% Lösungen Einheit 1 von 4 – Seite berechnen mit Sinus, Kosinus und Tangens (zusammenbau v0.1)
\einheitenkopf{Einheit 1 von 4 · Seite berechnen mit Sinus, Kosinus und Tangens}

\erg{12}{a) sin}
\erg{13}{a) im Zähler: mal}
%% TODO Lösung der Erklärzeile 14g) fehlt in der Bank
\erg{14}{a) r: G, s: A, t: H \quad b) $\mathrm{sin}\,\beta = \frac{e}{f}$ \quad c) $\mathrm{cos}\,\alpha = \frac{n}{m}$ \quad d) $\mathrm{tan}\,\gamma = \frac{z}{y}$ \quad e) $\mathrm{sin}\,\alpha = \frac{a}{c}$ \quad f) $\mathrm{cos}\,\beta = \frac{l}{m}$ \quad g) \ldots \quad h) $\mathrm{sin}\,\beta = \frac{q}{p}$ \quad i) $\mathrm{sin}\,23^\circ \approx 0{,}391$ \quad j) $\mathrm{sin}\,31^\circ = \frac{x}{8}$, $x = 8 \cdot \mathrm{sin}\,31^\circ \approx 4{,}1$ cm \quad k) $\mathrm{cos}\,27^\circ = \frac{x}{9}$, $x = 9 \cdot \mathrm{cos}\,27^\circ \approx 8{,}0$ cm \quad l) $b$ ist Gegenkathete von $\beta$: $b = 7 \cdot \mathrm{sin}\,41^\circ \approx 4{,}6$ cm \quad m) $\mathrm{sin}\,44^\circ = \frac{6}{x}$, $x = 6 : \mathrm{sin}\,44^\circ \approx 8{,}6$ cm \quad n) $\mathrm{cos}\,32^\circ = \frac{7}{x}$, $x = 7 : \mathrm{cos}\,32^\circ \approx 8{,}3$ cm \quad o) $\mathrm{tan}\,27^\circ = \frac{x}{6}$, $x = 6 \cdot \mathrm{tan}\,27^\circ \approx 3{,}1$ cm \quad p) $\mathrm{tan}\,31^\circ = \frac{5}{x}$, $x = 5 : \mathrm{tan}\,31^\circ \approx 8{,}3$ cm \quad q) Kosinus (Ankathete und Hypotenuse): $x = 9 \cdot \mathrm{cos}\,72^\circ \approx 2{,}8$ cm \quad r) $x = 6 : \mathrm{sin}\,48^\circ \approx 8{,}1$ \quad s) $x = 6{,}8 \cdot \mathrm{sin}\,42{,}6^\circ \approx 4{,}6$ m \quad t) $x = 7 \cdot \mathrm{cos}\,42^\circ \approx 5{,}20$ cm, also rund $5{,}2$ cm \quad u) Höhe $= 4{,}2 \cdot \mathrm{sin}\,71^\circ \approx 4{,}0$ m \quad v) $\mathrm{sin}\,47^\circ = \frac{AB}{840}$, $AB = 840 \cdot \mathrm{sin}\,47^\circ \approx 614$ m}
\erg{15}{a) $a$ ist jetzt Ankathete, $b$ Gegenkathete}
\erg{16}{a) $3{,}4 : 5 = 0{,}68$; $\mathrm{sin}\,43^\circ \approx 0{,}682$ – fast gleich}
\erg{17}{a) Nein – jede Kathete ist kürzer als die Hypotenuse.}
\erg{18}{a) Fehler: mal statt geteilt – $x$ steht im Nenner. Richtig: $x = 5 : \mathrm{sin}\,41^\circ \approx 7{,}6$ cm \quad b) Die Gegenkathete ist kürzer als die Hypotenuse, darum ist der Bruch Gegenkathete durch Hypotenuse kleiner als 1. \quad c) Höhe $= 6 \cdot \mathrm{sin}\,72^\circ \approx 5{,}7$ m – mehr als $5{,}5$ m, die Leiter reicht.}
